\section{引言}
协同组件的随机出现，是有限资源系统中的一类组合可及性问题。
单个组件进入样本的概率，无法直接说明多个组件何时共同出现；
即使组件齐全，后续收益也会随剩余资源构成而波动。
例如，一副24张牌的牌组中，某张指定牌进入前10张的概率为$10/24$，
而两个指定能力组件与一个可替代消耗组件共同进入前10张，需要计算另一个联合事件。
区分单组件出现、联合可及性与条件回报，是评价此类系统的基础。

《杀戮尖塔2》是一款以随机牌组构筑为核心的卡牌游戏\cite{steam}。
其中，消耗一张牌可以同时触发抽牌与防御收益，批量消耗则把一次随机手牌的类别构成转化为数值回报。
本文选择这一机制作为有限无放回抽样的实例：把能力组件、消耗组件与普通牌划分类别，
以固定观察预算研究组件形成的机会，再研究同一排列中随后样本的收益。
模型中的“格挡”表示可累加的防御点数。

本文围绕三个相互衔接的问题展开：关键组件的联合可及性与首次见齐时间如何变化，
同洗牌条件依赖如何决定消耗收益分布，以及构筑操作如何分别改变启动机会、条件回报和总体波动。
组合计数给出精确概率，交换性给出均值与方差关于牌类数量的结构表达，
全期望与全方差则把条件收益连接到总体指标。
本文以概率统计的基本工具\cite{blitzstein}进行推导，并用随机排列模拟检验数值结果及估计精度。

\clearpage
\section{抽象卡组与随机实验}\label{sec:model}
\subsection{类别、参数与样本空间}
设六类牌的数量分别为$a,b,c,d,t,s$，总数
$N=a+b+c+d+t+s$。基准构成见表\ref{tab:deck}。
同类牌仍有不同实体编号，样本空间由$N!$个等概率排列组成。
前$n$张是均匀的$n$元不放回样本；$h$表示其后的收益抽样张数，
满足$0\le n\le N$、$0\le h\le N-n$。
基准取$(N,n,h)=(24,10,4)$。
基准卡组包含核心能力、消耗组件和普通牌。
取$n=10$作为前期组件观察窗口，$h=4$作为收益抽样窗口，
使一张保留的C与四张随机牌组成五张收益手牌。
后文将比较不同卡组规模与类别数量下的概率与收益。

\begin{table}[htbp]
\centering\small
\caption{基准抽象卡组及固定机制}\label{tab:deck}
\begin{tabular}{clrp{9.3cm}}
\toprule
类别 & 机制原型 & 数量 & 模型中的作用\\
\midrule
A & 黑暗之拥 & 1 & 生效后，每消耗一张牌产生一次抽牌触发\\
B & 无惧疼痛 & 1 & 生效后，每消耗一张牌增加3点格挡\\
C & 重振精神 & 2 & 消耗其他非攻击牌，每张提供5点格挡\\
D & 恶魔之焰 & 1 & 消耗其余手牌；作为静态可及性分析中的备选消耗组件\\
T & 普通攻击牌 & 9 & 不属于C的消耗对象\\
S & 普通非攻击牌 & 10 & 可由C消耗\\
\bottomrule
\end{tabular}
\end{table}

A、B、C、S为非攻击类，D、T为攻击类。
公开卡牌资料中的抽牌、额外格挡、非攻击牌批量消耗与手牌批量消耗描述，
分别对应表中的四个机制原型\cite{dark,pain,wind,fire}。
本文采用未升级卡牌的基础效果，将每张消耗获得的$5+3=8$点格挡作为模型参数。

\subsection{可及性、顺序事件与两阶段协议}
定义$F_n^C$为前$n$张中A、B、C各至少出现一张的事件，
$F_n^E$则把C条件替换为“C或D至少出现一张”。
定义$O_n^C$、$O_n^E$为相应前缀中存在一张触发牌，
其位置严格晚于第一张A和第一张B的事件。
因此$O_n^q\subseteq F_n^q$，$q\in\{C,E\}$。
前缀较早出现的触发牌，可以与后来另一张满足顺序的触发牌并存。

收益实验使用同一排列，并依次执行以下协议。
揭示前$n$张，令$I=\ind_{F_n^C}$。
当$I=1$时，假想各激活一张A、B，保留一张已揭示C，其余已揭示牌退出本次实验；
收益手牌由保留的C与排列位置$n+1,\ldots,n+h$的牌组成。
令$X$为这$h$张中可被C消耗的非攻击牌数；当$I=0$时约定$X=0$。
C自身不计入$X$。条件收益为$G=8X$，总体收益定义为
\begin{equation}\label{eq:payoff-definition}
Y=IG=8IX,\qquad
p=\Prob(I=1),\quad
\mu=\E(G\mid I=1),\quad
\sigma^2=\V(G\mid I=1).
\end{equation}
此处条件均值与方差要求$p>0$。基准中A、B已经位于前缀，
故收益阶段的非攻击牌恰为C或S。

收益协议以$F_n^C$为成功条件，$O_n^C$用于单独考察出现顺序。
收益阶段统计A产生的抽牌触发次数$R=X$，并计算C与B共同产生的格挡；
D用于可及性比较。整个实验针对固定预算下的随机排列与消耗收益，
能量、抽牌反馈和牌堆循环作为后续的动态问题。
推广到多副本A/B时，协议各激活一张，其余未揭示副本按非攻击牌处理。
